% GENERATED FILE. Edit canonical lessons, then run npm run generate:books.
% canonical-source-hash: fnv1a64:e24631f75f7af4d4
% canonical-lessons: LA-C491-lacrimo, LA-C491-ploro, LA-C491-gemo, LA-C491-gemitus, LA-C491-suspiro

\chapter{Tears and sighs}
\label{ch:la-tears-and-sighs}
\hlchaptermodality{491}

\begin{chapteropening}
\textbf{By the end of this chapter:} \emph{I can describe weeping, groaning and sighing.}

Complete the last lesson of chapter 491: I can describe weeping, groaning and sighing.
\end{chapteropening}

\section[lacrimō, lacrimāre]{lacrimō, lacrimāre — to weep, shed tears}

\label{lesson:LA-C491-lacrimo}



\begin{quote}
Before the new one: say the Latin for grief, then the Latin for sorrowful.
\end{quote}

\subsection*{You'll want to know: lacrimō, lacrimāre}
\textbf{lacrimō, lacrimāre} — \textquotedblleft{}to weep, shed tears\textquotedblright{}.

\begin{quote}
Īnfāns tōtam noctem lacrimat.
\end{quote}

\begin{quote}
\emph{The baby cries all night.}
\end{quote}

\emph{lacrimō, lacrimāre, lacrimāvī, lacrimātum}, from \emph{lacrima} — a tear. English \emph{lachrymose} comes from the same family.

\begin{cousinweb}
\textbf{lacrimō} is built on \emph{lacrima} \textquotedblleft{}a tear\textquotedblright{} (Lewis \& Short), which EtymDB traces to a Proto-Indo-European word for \textquotedblleft{}tear\textquotedblright{}; English \emph{tear} is a cousin.
\end{cousinweb}

\begin{grammarlens}[title={What you've built}]
One more word for telling, explaining and describing at B1.
\end{grammarlens}

\subsection*{Your turn}
\begin{itemize}
\raggedright
  \item \emph{Say it:} \emph{lacrimō, lacrimāre}
  \item \emph{Say it:} \emph{lacrimō, lacrimāre}, once more
  \item \emph{Say it:} a sentence of your own with \emph{lacrimō, lacrimāre}
  \item \emph{Recall:} say the Latin for grief, then the Latin for sorrowful, then say \emph{lacrimō, lacrimāre} again
\end{itemize}

\begin{tcolorbox}[breakable,colback=teal!4,colframe=teal!35!black,title={Before you move on}]
What does lacrimō, lacrimāre mean? (\textquotedblleft{}To weep, shed tears\textquotedblright{}.) Say it once more.
\end{tcolorbox}
\section[plōrō, plōrāre]{plōrō, plōrāre — to wail, cry aloud}

\label{lesson:LA-C491-ploro}



\begin{quote}
Before the new one: say the Latin for sorrowful, then the Latin for to weep.
\end{quote}

\subsection*{You'll want to know: plōrō, plōrāre}
\textbf{plōrō, plōrāre} — \textquotedblleft{}to wail, cry aloud\textquotedblright{}.

\begin{quote}
Puella plōrat, quia pila eius in flūmen cecidit.
\end{quote}

\begin{quote}
\emph{The girl is wailing because her ball fell into the river.}
\end{quote}

\emph{plōrō, plōrāre, plōrāvī, plōrātum}. Louder than \emph{fleō}: crying out loud. English \emph{deplore} and \emph{implore} come from its compounds.

\begin{cousinweb}
\textbf{plōrō} is a native Latin verb whose origin is not settled; Lewis \& Short calls it doubtful. English \emph{deplore} and \emph{implore} come from its compounds \emph{dēplōrō} and \emph{implōrō}.
\end{cousinweb}

\begin{grammarlens}[title={What you've built}]
One more word for telling, explaining and describing at B1.
\end{grammarlens}

\subsection*{Your turn}
\begin{itemize}
\raggedright
  \item \emph{Say it:} \emph{plōrō, plōrāre}
  \item \emph{Say it:} \emph{plōrō, plōrāre}, once more
  \item \emph{Say it:} a sentence of your own with \emph{plōrō, plōrāre}
  \item \emph{Recall:} say the Latin for sorrowful, then the Latin for to weep, then say \emph{plōrō, plōrāre} again
\end{itemize}

\begin{tcolorbox}[breakable,colback=teal!4,colframe=teal!35!black,title={Before you move on}]
What does plōrō, plōrāre mean? (\textquotedblleft{}To wail, cry aloud\textquotedblright{}.) Say it once more.
\end{tcolorbox}
\section[gemō, gemere]{gemō, gemere — to groan, moan}

\label{lesson:LA-C491-gemo}



\begin{quote}
Before the new one: say the Latin for to weep, then the Latin for to wail.
\end{quote}

\subsection*{You'll want to know: gemō, gemere}
\textbf{gemō, gemere} — \textquotedblleft{}to groan, moan\textquotedblright{}.

\begin{quote}
Aeger in lectō gemit.
\end{quote}

\begin{quote}
\emph{The sick man groans in bed.}
\end{quote}

\emph{gemō, gemere, gemuī, gemitum}. Also of a dove cooing or a wheel creaking.

\begin{cousinweb}
\textbf{gemō} is a native Latin verb; Lewis \& Short compares a Greek verb meaning \textquotedblleft{}to be full\textquotedblright{}, but EtymDB gives no parent, so the origin is not settled.
\end{cousinweb}

\begin{grammarlens}[title={What you've built}]
One more word for telling, explaining and describing at B1.
\end{grammarlens}

\subsection*{Your turn}
\begin{itemize}
\raggedright
  \item \emph{Say it:} \emph{gemō, gemere}
  \item \emph{Say it:} \emph{gemō, gemere}, once more
  \item \emph{Say it:} a sentence of your own with \emph{gemō, gemere}
  \item \emph{Recall:} say the Latin for to weep, then the Latin for to wail, then say \emph{gemō, gemere} again
\end{itemize}

\begin{tcolorbox}[breakable,colback=teal!4,colframe=teal!35!black,title={Before you move on}]
What does gemō, gemere mean? (\textquotedblleft{}To groan, moan\textquotedblright{}.) Say it once more.
\end{tcolorbox}
\section[gemitus, gemitūs]{gemitus, gemitūs — a groan, moan}

\label{lesson:LA-C491-gemitus}



\begin{quote}
Before the new one: say the Latin for to wail, then the Latin for to groan.
\end{quote}

\subsection*{You'll want to know: gemitus, gemitūs}
\textbf{gemitus, gemitūs} — \textquotedblleft{}a groan, moan\textquotedblright{}.

\begin{quote}
Ē cubiculō gemitūs audīmus.
\end{quote}

\begin{quote}
\emph{We hear groans from the bedroom.}
\end{quote}

Masculine, fourth declension, from \emph{gemō}.

\begin{cousinweb}
\textbf{gemitus} is the action noun of \emph{gemō} \textquotedblleft{}to groan\textquotedblright{}, with the fourth-declension ending \emph{-tus} (Lewis \& Short; EtymDB derives it from \emph{gemō}).
\end{cousinweb}

\begin{grammarlens}[title={What you've built}]
One more word for telling, explaining and describing at B1.
\end{grammarlens}

\subsection*{Your turn}
\begin{itemize}
\raggedright
  \item \emph{Say it:} \emph{gemitus, gemitūs}
  \item \emph{Say it:} \emph{gemitus, gemitūs}, once more
  \item \emph{Say it:} a sentence of your own with \emph{gemitus, gemitūs}
  \item \emph{Recall:} say the Latin for to wail, then the Latin for to groan, then say \emph{gemitus, gemitūs} again
\end{itemize}

\begin{tcolorbox}[breakable,colback=teal!4,colframe=teal!35!black,title={Before you move on}]
What does gemitus, gemitūs mean? (\textquotedblleft{}A groan, moan\textquotedblright{}.) Say it once more.
\end{tcolorbox}
\section[suspīrō, suspīrāre]{suspīrō, suspīrāre — to sigh}

\label{lesson:LA-C491-suspiro}



\begin{quote}
Before the new one: say the Latin for to groan, then the Latin for a groan.
\end{quote}

\subsection*{You'll want to know: suspīrō, suspīrāre}
\textbf{suspīrō, suspīrāre} — \textquotedblleft{}to sigh\textquotedblright{}.

\begin{quote}
Iuvenis suspīrat et dē puellā cōgitat.
\end{quote}

\begin{quote}
\emph{The young man sighs and thinks about the girl.}
\end{quote}

\emph{suspīrō, suspīrāre, suspīrāvī, suspīrātum}: \emph{sub} 'up from below' + \emph{spīrō} 'to breathe'.

\begin{cousinweb}
\textbf{suspīrō} is built on \emph{sub-} \textquotedblleft{}up from below\textquotedblright{} and \emph{spīrō} \textquotedblleft{}to breathe\textquotedblright{} (Lewis \& Short; EtymDB): a sigh is a breath drawn up from deep down. English \emph{inspire} and \emph{respire} share \emph{spīrō}.
\end{cousinweb}

\begin{grammarlens}[title={What you've built}]
One more word for telling, explaining and describing at B1.
\end{grammarlens}

\subsection*{Your turn}
\begin{itemize}
\raggedright
  \item \emph{Say it:} \emph{suspīrō, suspīrāre}
  \item \emph{Say it:} \emph{suspīrō, suspīrāre}, once more
  \item \emph{Say it:} a sentence of your own with \emph{suspīrō, suspīrāre}
  \item \emph{Recall:} say the Latin for to groan, then the Latin for a groan, then say \emph{suspīrō, suspīrāre} again
\end{itemize}

\begin{tcolorbox}[breakable,colback=teal!4,colframe=teal!35!black,title={Before you move on}]
What does suspīrō, suspīrāre mean? (\textquotedblleft{}To sigh\textquotedblright{}.) Say it once more.
\end{tcolorbox}
